% GENERATED FILE. Edit canonical lessons, then run npm run generate:books.
% canonical-source-hash: fnv1a64:c3ec8fc8d877b30b
% canonical-lessons: IT-C247-bottone, IT-C247-taglia, IT-C247-colore, IT-C247-insalata, IT-C247-spuntino

\chapter{Button, Size, Colour, Salad, Snack}
\label{ch:it-button-size-colour-salad-snack}
\hlchaptermodality{247}

\begin{chapteropening}
\textbf{By the end of this chapter:} \emph{I can say a button, a size, a colour, salad and a snack.}

Complete the last lesson of chapter 247: I can say a button, a size, a colour, salad and a snack.
\end{chapteropening}

\section[il bottone]{il bottone — a button}

\label{lesson:IT-C247-bottone}



\begin{quote}
Before the new one: say the Italian for raw, then the Italian for cooked.
\end{quote}

\subsection*{You'll want to know: il bottone}
\textbf{il bottone} — \textquotedblleft{}a button\textquotedblright{}.

\begin{cousinweb}
\textbf{bottone} $\leftarrow$ Old French \emph{boton}, \textquotedblleft{}a bud, a button\textquotedblright{} — borrowed. EtymDB 2.1 has no row for \emph{bottone}; it builds the French word on \emph{bouter}, \textquotedblleft{}to push\textquotedblright{}.
\end{cousinweb}

\begin{grammarlens}[title={What you've built}]
One more word for this chapter.
\end{grammarlens}

\subsection*{Your turn}
\begin{itemize}
\raggedright
  \item \emph{Say it:} \emph{il bottone}
  \item \emph{Say it:} \emph{il bottone}, once more
  \item \emph{Say it:} say \emph{il bottone}
  \item \emph{Recall:} say the Italian for raw, then the Italian for cooked, then say \emph{il bottone} again
\end{itemize}

\begin{tcolorbox}[breakable,colback=teal!4,colframe=teal!35!black,title={Before you move on}]
What does il bottone mean? (\textquotedblleft{}A button\textquotedblright{}.) Say it once more.
\end{tcolorbox}
\section[la taglia]{la taglia — a size (of clothes)}

\label{lesson:IT-C247-taglia}



\begin{quote}
Before the new one: say the Italian for cooked, then the Italian for a button.
\end{quote}

\subsection*{You'll want to know: la taglia}
\textbf{la taglia} — \textquotedblleft{}a size (of clothes)\textquotedblright{}.

\begin{cousinweb}
\textbf{taglia} is the noun of \emph{tagliare}, \textquotedblleft{}to cut\textquotedblright{}: the cut of a garment. \emph{tagliare} $\leftarrow$ Late Latin \emph{taliare}. (EtymDB 2.1)
\end{cousinweb}

\begin{grammarlens}[title={What you've built}]
One more word for this chapter.
\end{grammarlens}

\subsection*{Your turn}
\begin{itemize}
\raggedright
  \item \emph{Say it:} \emph{la taglia}
  \item \emph{Say it:} \emph{la taglia}, once more
  \item \emph{Say it:} say \emph{la taglia}
  \item \emph{Recall:} say the Italian for cooked, then the Italian for a button, then say \emph{la taglia} again
\end{itemize}

\begin{tcolorbox}[breakable,colback=teal!4,colframe=teal!35!black,title={Before you move on}]
What does la taglia mean? (\textquotedblleft{}A size (of clothes)\textquotedblright{}.) Say it once more.
\end{tcolorbox}
\section[il colore]{il colore — a colour}

\label{lesson:IT-C247-colore}



\begin{quote}
Before the new one: say the Italian for a button, then the Italian for a size.
\end{quote}

\subsection*{You'll want to know: il colore}
\textbf{il colore} — \textquotedblleft{}a colour\textquotedblright{}.

\begin{cousinweb}
\textbf{colore} $\leftarrow$ Latin \emph{color}, \textquotedblleft{}colour\textquotedblright{} — inherited. (EtymDB 2.1)
\end{cousinweb}

\begin{grammarlens}[title={What you've built}]
One more word for this chapter.
\end{grammarlens}

\subsection*{Your turn}
\begin{itemize}
\raggedright
  \item \emph{Say it:} \emph{il colore}
  \item \emph{Say it:} \emph{il colore}, once more
  \item \emph{Say it:} say \emph{il colore}
  \item \emph{Recall:} say the Italian for a button, then the Italian for a size, then say \emph{il colore} again
\end{itemize}

\begin{tcolorbox}[breakable,colback=teal!4,colframe=teal!35!black,title={Before you move on}]
What does il colore mean? (\textquotedblleft{}A colour\textquotedblright{}.) Say it once more.
\end{tcolorbox}
\section[l'insalata]{l'insalata — salad}

\label{lesson:IT-C247-insalata}



\begin{quote}
Before the new one: say the Italian for a size, then the Italian for a colour.
\end{quote}

\subsection*{You'll want to know: l'insalata}
\textbf{l'insalata} — \textquotedblleft{}salad\textquotedblright{}.

\begin{cousinweb}
\textbf{insalata} is \textquotedblleft{}salted (greens)\textquotedblright{}: the past participle of \emph{insalare}, \textquotedblleft{}to salt\textquotedblright{}, built on \emph{sale}, \textquotedblleft{}salt\textquotedblright{}. EtymDB 2.1 has no row for it; the parts are plain to see.
\end{cousinweb}

\begin{grammarlens}[title={What you've built}]
One more word for this chapter.
\end{grammarlens}

\subsection*{Your turn}
\begin{itemize}
\raggedright
  \item \emph{Say it:} \emph{l'insalata}
  \item \emph{Say it:} \emph{l'insalata}, once more
  \item \emph{Say it:} say \emph{l'insalata}
  \item \emph{Recall:} say the Italian for a size, then the Italian for a colour, then say \emph{l'insalata} again
\end{itemize}

\begin{tcolorbox}[breakable,colback=teal!4,colframe=teal!35!black,title={Before you move on}]
What does l'insalata mean? (\textquotedblleft{}Salad\textquotedblright{}.) Say it once more.
\end{tcolorbox}
\section[lo spuntino]{lo spuntino — a snack}

\label{lesson:IT-C247-spuntino}



\begin{quote}
Before the new one: say the Italian for a colour, then the Italian for salad.
\end{quote}

\subsection*{You'll want to know: lo spuntino}
\textbf{lo spuntino} — \textquotedblleft{}a snack\textquotedblright{}.

\begin{cousinweb}
\textbf{spuntino} — built from \emph{spuntare}, \textquotedblleft{}to take the point off\textquotedblright{}, with \emph{-ino}: a bite that blunts your hunger. EtymDB 2.1 has no row for it.
\end{cousinweb}

\begin{grammarlens}[title={What you've built}]
One more word for this chapter.
\end{grammarlens}

\subsection*{Your turn}
\begin{itemize}
\raggedright
  \item \emph{Say it:} \emph{lo spuntino}
  \item \emph{Say it:} \emph{lo spuntino}, once more
  \item \emph{Say it:} say \emph{lo spuntino}
  \item \emph{Recall:} say the Italian for a colour, then the Italian for salad, then say \emph{lo spuntino} again
\end{itemize}

\begin{tcolorbox}[breakable,colback=teal!4,colframe=teal!35!black,title={Before you move on}]
What does lo spuntino mean? (\textquotedblleft{}A snack\textquotedblright{}.) Say it once more.
\end{tcolorbox}
